% Figure: the six-bay Pratt truss of the worked example, with bays
% labelled, the section cut through the third bay shown as a dashed line,
% and the joint loads drawn as downward arrows on the bottom chord.
\begin{figure}[htbp]
\centering
\begin{tikzpicture}[>=latex, x=1.0cm, y=1.0cm]
  % bottom chord joints L0..L6 (12 m span, 6 bays of 2 m -> scale 1cm = 1m)
  \foreach \i in {0,...,6} {
    \coordinate (L\i) at (\i*2, 0);
  }
  % top chord joints U1..U5 above interior bottom joints (height 2 m)
  \foreach \i in {1,...,5} {
    \coordinate (U\i) at (\i*2, 2);
  }
  % bottom chord
  \draw[engblue, very thick] (L0) -- (L6);
  % top chord
  \draw[engblue, very thick] (U1) -- (U5);
  % end posts / inclined ends
  \draw[engblue, very thick] (L0) -- (U1);
  \draw[engblue, very thick] (L6) -- (U5);
  % verticals
  \foreach \i in {1,...,5} { \draw[engblue, thick] (L\i) -- (U\i); }
  % Pratt diagonals slope toward midspan (tension diagonals)
  \draw[engblue, thick] (U1) -- (L2);
  \draw[engblue, thick] (U2) -- (L3);
  \draw[engblue, thick] (L3) -- (U4);
  \draw[engblue, thick] (L4) -- (U5);
  % joints
  \foreach \i in {0,...,6} { \fill[engblack] (L\i) circle (1.6pt); }
  \foreach \i in {1,...,5} { \fill[engblack] (U\i) circle (1.6pt); }
  % joint loads at interior bottom joints
  \foreach \i in {1,...,5} {
    \draw[->, very thick, engred] (L\i) -- ++(0,-1.0);
    \node[engred, font=\scriptsize, anchor=north] at ($(L\i)+(0,-1.0)$) {$P$};
  }
  % supports
  \draw[thick] ($(L0)+(-0.3,-0.45)$) -- (L0) -- ($(L0)+(0.3,-0.45)$) -- cycle;
  \draw[thick] ($(L6)+(0,-0.18)$) circle (0.12);
  \draw[thick] ($(L6)+(-0.3,-0.3)$) -- ($(L6)+(0.3,-0.3)$);
  % reactions
  \draw[->, thick, englime!70!black] ($(L0)+(0,-0.55)$) -- ++(0,-0.7);
  \node[font=\scriptsize, anchor=east] at ($(L0)+(-0.05,-1.0)$) {$V_A$};
  \draw[->, thick, englime!70!black] ($(L6)+(0,-0.45)$) -- ++(0,-0.7);
  \node[font=\scriptsize, anchor=west] at ($(L6)+(0.05,-0.9)$) {$V_B$};
  % section cut through third bay (between L2 and L3): cuts top chord U2-U3?
  % the section just right of L2 cuts top chord U2-U3, diagonal U2-L3, bottom L2-L3
  \draw[engamber, dashed, very thick] (5,2.6) -- (5,-0.7);
  \node[engamber, font=\scriptsize, anchor=south] at (5,2.6) {section};
  % labels for a few joints
  \node[font=\scriptsize, anchor=north east] at (L0) {$L_0$};
  \node[font=\scriptsize, anchor=north] at ($(L3)+(0,0.0)$) {};
  \node[font=\scriptsize, anchor=south] at ($(U2)+(0,0.05)$) {$U_2$};
  \node[font=\scriptsize, anchor=south] at ($(U3)+(0,0.05)$) {$U_3$};
\end{tikzpicture}
\caption[Six-bay Pratt truss of the worked example]{The Pratt truss of
the worked example: span $L=12\,\si{\meter}$, depth $h=2\,\si{\meter}$,
six bays of $2\,\si{\meter}$, with equal downward loads $P=20\,\text{kN}$
at the five interior bottom-chord joints. The Pratt configuration runs
its diagonals down toward midspan so that under gravity load the
diagonals carry tension and the shorter verticals carry compression, the
arrangement that lets a designer use slender tension diagonals and stout
verticals. The amber dashed line is the section cut used to extract the
third-bay diagonal force in one calculation: it crosses exactly three
members (top chord $U_2U_3$, diagonal $U_2L_3$, bottom chord $L_2L_3$),
the maximum a planar section may cross and still be solvable by the three
planar equilibrium equations.}
\label{fig:vol03ch02:pratt-truss}
\end{figure}
